\chapter{Experimental Method}\label{ch:method}

The final experiment measures one specific use case: the system has already
seen earlier MLII ECG from a patient, and it must classify later ECG from that
same patient. This is a known-patient temporal holdout. It is not a
new-patient external validation and it is not a shuffled random beat split.

\section{Method overview}

Figure~\ref{fig:method-workflow} shows how accepted beats are allocated to
model selection, final refitting, and testing. The normal
autoencoder was already frozen from NSRDB. The supervised experiment then used
only MITDB MLII beats. Candidate models were selected with earlier temporal
regions, refitted before the test boundary, and evaluated only on the final
20\% of accepted beats.

\begin{figure}[H]
\centering
\includegraphics[width=.98\textwidth]{pic/generated/method_final_workflow.pdf}
\caption[Chronological model selection, refitting, and final testing.]{Chronological
selection, refitting, and testing by accepted-beat order. Selection excludes
16 beats before both boundaries; refitting retains only the exclusion before
the test. The test region was observed during earlier development.}
\label{fig:method-workflow}
\end{figure}

\section{Research objective and success conditions}

The final experiment asks:
\begin{quote}
After the model has access to the earlier ECG of each patient, how accurately
does it classify that patient's later ECG?
\end{quote}

The primary outcome is whole-system six-class accuracy on the final 20\%
temporal holdout. The planned success conditions were:
\begin{enumerate}
    \item pooled whole-system accuracy at least 90\%;
    \item lower bound of the equal-subject 95\% confidence interval above
    90\%; and
    \item total confidence-interval width no greater than three percentage
    points.
\end{enumerate}
Per-class recall and macro F1 were also checked so that a high overall score
could not hide failure on a smaller arrhythmia class.

\section{Input data and feature construction}

The method uses one ECG channel at every stage, but the two data sources have
different roles. NSRDB supplies healthy windows for the frozen normal
autoencoder. MITDB supplies MLII beat windows, binary labels, six-class labels,
and the final temporal test. No NSRDB beat is scored in the final test.

\begin{table}[H]
\centering
\caption{Exact evidence supplied to the final supervised models.}
\label{tab:method-feature-input}
\begin{tabular}{>{\raggedright\arraybackslash}p{.26\textwidth}
                >{\raggedright\arraybackslash}p{.20\textwidth}
                >{\raggedright\arraybackslash}p{.44\textwidth}}
\toprule
Evidence block & Count & Role\\
\midrule
MLII morphology and spectrum & 223 values & Describes the beat-centred
512-sample MLII waveform. This is the direct ECG evidence for both supervised
stages.\\
RR timing context & 15 values & Describes previous RR, next RR, local timing,
and recent rhythm history. In live use the beat is scored only after this
next-RR context is available; labels are never used during prediction.\\
Frozen autoencoder residuals & 4 values & Summarises how the frozen
NSRDB-trained normal autoencoder reconstructs the same MLII beat.\\
\midrule
Total per MITDB beat & 242 values & The binary detector and subtype classifier
receive the same one-lead feature vector.\\
\bottomrule
\end{tabular}
\end{table}

This feature construction matches the saved final model bundle:
\texttt{lead\_count} is one, \texttt{required\_lead} is MLII,
\texttt{rr\_feature\_count} is 15, and \texttt{ae\_feature\_count} is four.
The method therefore does not use a second MITDB lead during training,
validation, temporal testing, or deployment.

\section{Normal-reference autoencoder training procedure}

The autoencoder was trained \emph{before} either supervised model. All 18
NSRDB records supplied first-channel ECG1 data. After filtering,
four-second windowing, per-window standardisation, and quality checks,
15 entire subjects contributed 288,673 fitting windows and three different
subjects contributed 57,177 validation windows. The split is by subject, so
adjacent windows from the same recording cannot appear on opposite sides of
this validation boundary.

One training step takes a shuffled batch of 256 clean windows, adds fresh
Gaussian noise with standard deviation 0.03 to the inputs, reconstructs them,
and computes mean squared error against the original clean windows. PyTorch
backpropagates this error through the convolutional encoder and decoder;
Adam updates their weights with learning rate 0.001. The training code repeats
these steps for two complete epochs. After each epoch it disables dropout and
measures clean-input MSE on the three untouched validation subjects. It saves
the checkpoint with the lowest validation MSE. The precise objective and
gradient logic are in Equations~\ref{eq:theory-denoising-loss} and
\ref{eq:theory-adam}; Figure~\ref{fig:theory-training-mechanisms} contrasts
this loop with tree fitting. The measured loss curve appears in
Figure~\ref{fig:lead-ii-ae-learning}.

After checkpoint selection, the model is frozen. For every MITDB MLII window,
its reconstruction is compared with the input to create four residual
measurements. No MITDB class label updates the autoencoder. The saved
normal-residual threshold is descriptive for the normal-reference model;
the final classifier uses residual measurements as features and uses the
separately validated binary-tree threshold for its final gate.

\section{Temporal split}

The eligible supervised records are the 46 MITDB records that contain MLII.
Records 102 and 104 were excluded before fitting because their available leads
are V5 and V2, not MLII. Records 201 and 202 are treated as one subject cluster
for uncertainty analysis because they belong to the same person.

Every eligible record is kept in chronological order. Percentages refer to
accepted-beat counts, not exact elapsed recording time. With $n_r$ accepted
beats, the cuts are $\lfloor0.64n_r\rfloor$ and $\lfloor0.80n_r\rfloor$:
\begin{enumerate}
    \item the first 64\% is the candidate-training region;
    \item a 16-beat gap is removed before validation during candidate
    selection;
    \item the 64--80\% region is the validation region;
    \item a 16-beat gap is removed before the final temporal test; and
    \item the final 20\% is the fixed temporal test.
\end{enumerate}
For final fitting, the selected binary and subtype models are refitted on the
earlier 80\% development region after the 16-beat boundary gap before the test
is removed. The test segment itself remains intact.

\begin{figure}[H]
\centering
\includegraphics[width=.98\textwidth,height=.70\textheight,keepaspectratio]{pic/generated/method_temporal_split_audit.pdf}
\caption{Temporal split audit for every eligible MITDB MLII record. Blue beats
are available for final fitting, yellow beats are the boundary gap before the
test, and red beats are the held-out final 20\% temporal test.}
\label{fig:method-split-audit}
\end{figure}

The final held-out test contains 20,969 accepted beats. These are later beats
from known patients, not beats from unseen patients. The method prevents signal
overlap across the test boundary, but it intentionally keeps patient overlap.

\section{Model selection and final fitting}

Model selection uses only the earlier temporal validation region. In this
rerun, the final 20\% test region does not choose the candidate, threshold, or
hyperparameters.

Three binary Extra Trees candidates were fitted on the earlier candidate
region and ranked by validation balanced accuracy, then by F1. Three subtype
Extra Trees candidates were fitted only on supported abnormal candidate beats
and ranked by validation macro F1, then by accuracy. Figure~\ref{fig:method-model-selection}
shows the validation scores that selected the final pair.

The fitting sequence for each candidate is explicit:
\begin{enumerate}
    \item Build one 242-value row per accepted MLII beat, keeping the label and
    subject identifier outside the row. The autoencoder is already frozen.
    \item Select candidate-training rows from the first 64\% of each record,
    subject to the 16-beat gap before validation. For the binary stage, use
    every normal/abnormal target. For the subtype stage, keep only supported
    abnormal targets.
    \item Compute class-and-record sample weights within these selected rows.
    At each tree node, propose random feature thresholds and choose the one
    giving the largest weighted entropy decrease in
    Equation~\ref{eq:theory-information-gain}. Continue until that candidate's
    depth and leaf-size rules stop growth, then average all tree outputs.
    \item Score the candidate on the 64--80\% validation region. For each
    binary candidate, choose an abnormal-score threshold using validation
    balanced accuracy and tie-breaks; compare candidates by balanced accuracy
    then F1. Compare subtype candidates by macro F1 then accuracy.
    \item Fix the winning structures and threshold. Refit the binary model on
    all 83,050 eligible development beats and the subtype model on 28,758
    supported abnormal development beats. Apply the frozen pair once to the
    final temporal segment.
\end{enumerate}
Tree fitting is a recursive partitioning procedure, so the supervised models
have no epoch-by-epoch training-loss curve. Their training criterion is node
entropy; the evidence used for model selection is held-out validation
performance. The two kinds of decrease should not be conflated.

\begin{figure}[H]
\centering
\includegraphics[width=.98\textwidth]{pic/generated/method_model_selection.pdf}
\caption{Model selection on the 64--80\% validation region. The binary detector
is selected by balanced accuracy and F1; the subtype classifier is selected by
macro F1 and accuracy.}
\label{fig:method-model-selection}
\end{figure}

The selected binary detector was the fully grown 240-tree candidate with
square-root feature sampling, minimum leaf size one, and record-weighting
power 0.35. Its validation-selected abnormal threshold was 0.5708333. The
selected subtype classifier was the fully grown 260-tree candidate with
one-third feature sampling, minimum leaf size one, and record-weighting power
0.25. After selection, the binary model was refitted on all development beats,
whereas the subtype model was refitted only on supported abnormal development
beats.

\section{Evaluation order}

The deployed decision is hierarchical. First, the binary detector estimates
whether a beat is abnormal. If the abnormal probability is not above the fixed
threshold, the final output is N. If the beat passes the gate, the subtype
classifier chooses PVC, PAC, LBBB, RBBB, or AFib.

\begin{figure}[H]
\centering
\includegraphics[width=.98\textwidth]{pic/generated/method_evaluation_order.pdf}
\caption[Populations used for the three evaluation questions.]{Populations used
for the three evaluation questions. Conditional subtype scoring uses true
supported abnormal beats independently of the gate. Complete-system scoring
uses the gate and subtype model together. The populations overlap; unsupported
abnormal labels enter only the binary evaluation.}
\label{fig:method-evaluation-order}
\end{figure}

The temporal evaluation reports three supports:
\begin{enumerate}
    \item binary abnormality detection on all 20,969 held-out beats, including
    unsupported abnormal MITDB symbols;
    \item conditional subtype classification on 7,958 true supported abnormal
    beats; and
    \item whole-system N/PVC/PAC/LBBB/RBBB/AFib classification on 20,043
    supported held-out beats.
\end{enumerate}
The 926 unsupported abnormal beats are useful for testing the binary
normal/abnormal gate, but they are not scored in the six-class whole-system
confusion matrix.

\section{Metrics}

For a binary outcome,
\begin{align}
 \mathrm{Accuracy}&=\frac{TP+TN}{TP+TN+FP+FN},\\
 \mathrm{Sensitivity}&=\frac{TP}{TP+FN},\\
 \mathrm{Specificity}&=\frac{TN}{TN+FP},\\
 \mathrm{Balanced\ accuracy}&=\frac{\mathrm{Sensitivity}+\mathrm{Specificity}}{2}.
\end{align}
For class $c$,
\begin{align}
 \mathrm{Precision}_c&=\frac{TP_c}{TP_c+FP_c},\\
 \mathrm{Recall}_c&=\frac{TP_c}{TP_c+FN_c},\\
 F1_c&=\frac{2\,\mathrm{Precision}_c\mathrm{Recall}_c}
 {\mathrm{Precision}_c+\mathrm{Recall}_c}.
\end{align}
Macro F1 is the unweighted mean of class F1 values. It gives PAC the same
weight as N even though N contains many more beats.

Two accuracy summaries are reported. Pooled beat accuracy gives every beat
equal weight. Subject-macro accuracy first calculates accuracy inside each of
the 45 eligible subject clusters and then gives every subject equal weight:
\begin{equation}
 A_{macro}=\frac{1}{45}\sum_{s=1}^{45}A_s.
 \label{eq:subject-macro}
\end{equation}
The subject-macro value is primary for uncertainty because a long record should
not behave like many independent patients.

\section{Subject-clustered confidence interval}

The primary 95\% confidence interval uses a non-parametric percentile
bootstrap:
\begin{enumerate}
    \item keep every beat from one subject together;
    \item treat records 201 and 202 as one subject cluster;
    \item sample 45 subject clusters with replacement;
    \item calculate the equal-subject mean accuracy;
    \item repeat 10,000 times using seed 20260801; and
    \item take the 2.5th and 97.5th percentiles.
\end{enumerate}
These intervals describe variation across the evaluated subject clusters with
the fitted model held fixed. They do not include retraining variability,
uncertainty from earlier test exposure, or a new population shift.
Resampling individual beats would create an artificially narrow interval
because adjacent beats from the same recording are correlated. The clustered
bootstrap keeps that dependence \cite{rutter2000clustered}.

\section{Tests against the 90\% target}

The confidence interval is the primary statistical evidence. Two
non-parametric tests provide sensitivity analyses. The exact one-sided sign
test asks whether more than half of subjects exceed 90\%. The one-sided
Wilcoxon signed-rank test asks whether subject accuracies are shifted above
90\%, under the additional symmetry assumption of a signed-rank location
test. Accuracy is bounded and has ceiling effects, so the Wilcoxon analysis is
secondary. The sign test concerns the fraction of subjects above the target,
not pooled beat accuracy. These tests are interpreted together with the confidence interval,
interval width, per-class recall, macro F1, and subject accuracy distribution.

\section{Reproducibility and leakage controls}

The final experiment saves the temporal boundary audit for every record, the
candidate scores, selected models, probability threshold, final confusion
matrix, per-class counts, per-subject accuracies, bootstrap seed, and headline
metrics. The rerun used seed 20260801 and did not use the final 20\% segment
for candidate selection, threshold selection, or final refitting.

The final 20\% segment had been observed during earlier system development, so
the evaluation is a retrospective rerun on previously observed data.
This prior exposure can influence design choices and cause optimism even
when the present selection code excludes test rows. It is not a confirmatory
first-use test, and the bootstrap does not correct development-selection bias. Patient overlap is
also intentional: the test is later ECG from known patients. Therefore, the
method supports a patient-adapted temporal-performance claim, not a
patient-independent generalisation claim.
